\chapter*{Introduction générale}
\addcontentsline{toc}{chapter}{Introduction générale}

Les essais cliniques produisent de nombreux documents : comptes rendus de biologie, d'imagerie, courriers et notes de suivi. Une partie des informations qu'ils contiennent doit être recopiée dans des formulaires électroniques de recueil de données, appelés eCRF. Cette opération est indispensable à la qualité des données d'étude, mais elle est longue, répétitive et exposée aux erreurs de transcription. Elle est également rendue difficile par la diversité des mises en page, des unités et des vocabulaires employés par les établissements.

Dans ce contexte, le stage réalisé chez AMALYTICS a porté sur le développement d'une solution d'extraction automatique d'informations à partir de documents médicaux. Le but n'est pas de prendre une décision clinique, ni de remplacer la validation humaine. Il s'agit d'assister le flux de saisie en proposant des valeurs structurées, accompagnées de leur provenance et soumises à des règles de contrôle avant export.

Le projet, nommé \emph{rapid-flow} dans le dépôt de développement, privilégie une architecture modulaire. Chaque étape traite une responsabilité précise : lecture du document, identification de son type, segmentation, recherche de contexte, extraction adaptée à une famille de champs, normalisation, validation et génération d'un export. Cette granularité permet de tester et de faire évoluer la chaîne sans confondre les causes d'une erreur.

Les travaux réalisés pendant le stage ont visé quatre objectifs : fiabiliser l'analyse des PDF, mettre en place un découpage qui préserve le contexte médical, renforcer l'extraction des bilans biologiques et de l'imagerie, puis fournir une évaluation reproductible. L'implémentation intègre aussi une recherche hybride, une isolation par étude et des politiques d'écriture qui évitent l'écrasement accidentel de données existantes.

Ce rapport est organisé en quatre chapitres. Le premier présente le contexte, l'étude de l'existant et les besoins. Le deuxième expose l'analyse et la conception. Le troisième décrit la réalisation de la solution. Le quatrième détaille la validation, les résultats obtenus et leurs limites. La conclusion récapitule les apports du stage et les perspectives envisageables.
